\section{Conclusiones}

El método de Euler permitió caracterizar el comportamiento transitorio y el
régimen estacional del modelo. Para el caso nominal, la precipitación presenta
una severidad estacional de aproximadamente $0.43$ y su máximo ocurre con un
retraso cercano a $0.8$ meses respecto del máximo de temperatura.

El análisis de sensibilidad muestra que los parámetros no actúan de manera
independiente: $A$ controla la intensidad de la excitación estacional,
$\beta$ favorece la generación de humedad, $\gamma$ regula su consumo por la
precipitación y $\sigma$ la disipación de esta última. La combinación
$(A,\beta,\gamma,\sigma)=(0.60,0.30,0.05,1.00)$ produjo la máxima severidad
explorada, $S_P=0.9929$, manteniendo un régimen anual repetitivo. En
particular, que $\gamma$ tome su valor mínimo confirma que maximizar todos los
parámetros no conduce necesariamente al máximo contraste estacional.
